\documentclass[11pt]{amsart}

\usepackage{amsmath,amssymb,amsthm}
\usepackage[margin=2.6cm]{geometry}
\usepackage{booktabs}
\usepackage{longtable}
\usepackage{seqsplit}
\usepackage[hidelinks]{hyperref}

\newtheorem{theorem}{Theorem}[section]
\newtheorem{lemma}[theorem]{Lemma}
\newtheorem{proposition}[theorem]{Proposition}
\newtheorem{corollary}[theorem]{Corollary}
\theoremstyle{definition}
\newtheorem{definition}[theorem]{Definition}
\theoremstyle{remark}
\newtheorem{remark}[theorem]{Remark}

\newcommand{\lean}[1]{\texttt{\seqsplit{#1}}}
\newcommand{\Mov}{\mathrel{\sim}}
\newcommand{\Z}{\mathbb{Z}}
\DeclareMathOperator{\pos}{pos}
\DeclareMathOperator{\len}{len}

\title[J11: torus word grids]{J11: the grid of the torus word and the reflected torus grid}
\date{2 November 2026}

\begin{document}

\begin{abstract}
Let $a\ge2$ and $q\ge1$. Write $\delta=\sigma_0\sigma_1\cdots\sigma_{a-2}$ for the Coxeter word on
$a$ strands. We prove that the selected grid diagram of the positive word $\delta^q$ and the row
reflection of TauCeti's grid $\lean{torusLink (a-1) (q-1)}$ are related by grid moves. The case
$a=S$, $q=S+1$ gives
\[
\begin{aligned}
&\lean{reflectedTorusGrid (S-1)}\\
&\qquad\Mov\lean{wordGrid (fullTwistWord S ++ coxeterWord S)}.
\end{aligned}
\]
The proof replaces $\delta^q$ by the conjugate word $\beta^q$, where
$\beta=\sigma_{a-2}\cdots\sigma_1\sigma_0$, and absorbs the factors $\beta$ one at a time into a
growing torus grid by explicit commutations and destabilizations. A preliminary computation of
Jones polynomials and linking numbers confirms that no chirality obstruction occurs.
\end{abstract}

\maketitle

\section{Statement, conventions and sources}\label{sec:statement}

\subsection{Grid diagrams}
We use the vendored TauCeti conventions. A grid diagram of size $n$ is a pair $(O,X)$ of
permutations of $\{0,\dots,n-1\}$, read as maps from columns to rows, with $O(c)\ne X(c)$ for every
column $c$ (\lean{TauCeti.GridDiagram}). For $x\ne y$ in $\{0,\dots,n-1\}$ the open cyclic
interval $(x,y)^{\circ}$ is $\{z: x<z<y\}$ if $x<y$ and $\{z: z>x\}\cup\{z: z<y\}$ if $x>y$
(\lean{Grid.cIoo}). Two pairs $\{x_0,x_1\}$, $\{y_0,y_1\}$ are \emph{non-interleaving} if
$x_0\in(y_0,y_1)^\circ\Leftrightarrow x_1\in(y_0,y_1)^\circ$ and
$y_0\in(x_0,x_1)^\circ\Leftrightarrow y_1\in(x_0,x_1)^\circ$ (\lean{Grid.Noninterleaving}).

The elementary moves are those of \lean{GridDiagram.IsMove} in
\lean{External/TauCeti/KnotTheory/Grid/Move/Basic.lean}:
\begin{enumerate}
\item a cyclic permutation of all rows or of all columns (\lean{relabelRows (finRotate n)},
\lean{relabelColumns (finRotate n)});
\item the exchange of two cyclically consecutive columns $c$, $c+1$ whose segments
$\{O(c),X(c)\}$ and $\{O(c+1),X(c+1)\}$ are non-interleaving, and the analogous exchange of two
consecutive rows, where the segment of a row is the pair of columns of its two markings
(\lean{movesTo\_swapColumns}, \lean{movesTo\_swapRows});
\item a stabilization or destabilization of one of the eight types (\lean{IsStabilization},
\lean{IsDestabilization}).
\end{enumerate}
We write $G\Mov G'$ for the relation \lean{MovesTo}, the equivalence relation generated by these
moves. It relates diagrams of different sizes.

For destabilizations we use the following description, which is the content of
\lean{isXStabilization\_eraseOCorner} in \lean{GridMoveRecognition.lean} and of its
mirror image obtained by exchanging $O$ and $X$ (\lean{isDestabilization\_swapMarkings}).
Suppose that columns $c,d$ are cyclically consecutive and rows $r,s$ are cyclically
consecutive, and that the corners $(c,r)$, $(d,r)$, $(c,s)$ carry the markings $O,X,X$ while
$(d,s)$ carries no $O$. Deleting column $c$ and row $r$, and moving the $X$-marking of column
$d$ to row $s$, is an $X$-destabilization. Exchanging the letters $O$ and $X$ in this sentence
gives an $O$-destabilization.

\subsection{The statement}
The Lean definitions involved are the following.
\begin{itemize}
\item \lean{coxeterWord a = List.finRange (a-1)}, the word $\delta$;
\lean{torusWord a q} $=\delta^q$, and \lean{fullTwistWord a} $=\delta^a$
(\lean{PositiveWordNormalForms.lean}). By \lean{torusWord\_succ},
$\delta^{S+1}=\delta^{S}\delta$, so \lean{torusWord S (S+1) = fullTwistWord S ++ coxeterWord S}.
\item \lean{wordGrid w}, the grid field of the selected data of a word $w$. It is an admissible
grid of $w$ in the sense of Section~\ref{sec:labeled} (\lean{wordGrid\_isAdmissible} in
\lean{PositiveWordGridAdmissible.lean}).
\item \lean{GridDiagram.torusLink p q}, of size $N=(p+1)+(q+1)$, with $O(c)=c$ and $X(c)=c+q+1$
modulo $N$ (\lean{torusLink\_O\_apply}, \lean{torusLink\_X\_apply}).
\item \lean{reflectedDiagram G = G.relabelRows Fin.revPerm}, so the reflected diagram has
$O(c)=N-1-O_G(c)$ and $X(c)=N-1-X_G(c)$; and
\lean{reflectedTorusGrid n = reflectedDiagram (torusLink n (n+1))}.
\end{itemize}

\begin{theorem}\label{thm:main}
Let $a\ge2$ and $q\ge1$. Then
\[
\lean{wordGrid (torusWord a q)}\ \Mov\ \lean{reflectedDiagram (torusLink (a-1) (q-1))}.
\]
In particular, for $S\ge2$,
\[
\lean{reflectedTorusGrid (S-1)}\Mov\lean{wordGrid (fullTwistWord S ++ coxeterWord S)}.
\]
\end{theorem}

The second assertion is the case $a=S$, $q=S+1$ of the first, because
$\lean{torusLink (S-1) S}$ is the diagram reflected in \lean{reflectedTorusGrid (S-1)}.

\subsection{Sources}
The proof is self-contained apart from the cited Lean declarations. The admissible grids of a
word are described in chapter J10 of \lean{docs/route-n/written-proofs-J05-J06-J07-J21-J10-J15-J14.tex}
(section \emph{The row-separated word grids}) and in \lean{PositiveWordGridAdmissible.lean}. The
statement that any two admissible grids of one word are related by row commutations is
\lean{IsAdmissibleWordGrid.movesTo} there. The conjugacy of Coxeter words under far commutations and
rotations is \lean{farRotationEquiv\_coxeterWord} of \lean{PositiveWordCoxeterConjugacy.lean} (job
J08). The vendored TauCeti file on torus grids cites Ozsv\'ath--Stipsicz--Szab\'o, \emph{Grid homology
for knots and links}, Chapter~3, which was not consulted. Nothing below depends on it.

\section{A chirality check}\label{sec:chirality}

The audit asks for a check that the two diagrams of Theorem~\ref{thm:main} have the same chirality.
We record a computation. It is evidence only, and the proof in later sections does not use it.

A grid diagram determines a planar diagram once one fixes which segment passes over the other at
a crossing. Any fixed choice gives a combinatorial invariant of the relation $\Mov$ only up to
mirror image, but it allows us to compare two diagrams. With vertical segments over horizontal
ones and the orientation from $X$ to $O$ in each column, we computed for both diagrams the
Kauffman bracket polynomial normalized by the writhe, and, for links, the pairwise linking
numbers. The computation (a Python script in the scratch directory of this job) gave equal values
for $\lean{wordGrid}(\delta^q)$ and $\lean{reflectedDiagram (torusLink (a-1) (q-1))}$ for
$(a,q)\in\{(2,2),(2,3),(2,4),(3,3),(2,5),(3,4),(3,2),(4,5)\}$, and different values for the
unreflected diagram $\lean{torusLink (a-1) (q-1)}$. All linking numbers of the word grids are
positive. For $(a,q)=(2,3)$ the normalized bracket is $A^{-4}+A^{-12}-A^{-16}$ for the word grid
and for the reflected torus grid. Thus the reflection in the statement is the expected one. A
greedy search for destabilizations and commutations also reduced each of these word grids to a
cyclic permutation of the reflected torus grid. The explicit sequence in
Section~\ref{sec:absorb} was found by analysing such reductions, and it was checked by machine for
$2\le a\le6$, $1\le k\le 5$ and for the words $v=\beta^j$, $0\le j\le2$. These checks are not part
of the proof.

\section{Labeled grid diagrams and elementary lemmas}\label{sec:labeled}

\begin{definition}\label{def:labeled}
A \emph{labeled grid diagram} consists of a sequence $C=(c_0,\dots,c_{n-1})$ of distinct column
labels, a sequence $R=(r_0,\dots,r_{n-1})$ of distinct row labels, and two maps $O,X$ from
$\{c_0,\dots,c_{n-1}\}$ to $\{r_0,\dots,r_{n-1}\}$ which are bijections and satisfy $O(c)\ne X(c)$ for
every column label $c$. Its \emph{position maps} are $\pos(c_i)=i$ and $\pos(r_j)=j$. The associated
grid diagram has $O'(i)=\pos(O(c_i))$ and $X'(i)=\pos(X(c_i))$.
\end{definition}

The associated grid diagram is a grid diagram because $O,X$ are bijections and $O(c)\ne X(c)$.
Relabeling the labels by bijections that preserve both sequences does not change it. We say that
a row label $r$ has its markings in the columns $O^{-1}(r)$ and $X^{-1}(r)$.

Moves are described on labeled diagrams. Exchanging two consecutive entries of $C$ (or of $R$)
produces a labeled diagram whose associated diagram is obtained by \lean{swapColumns} (or
\lean{swapRows}) of the two positions. The exchange is a grid move when the two segments are
non-interleaving. Moving the first entry of $C$ to the end is a cyclic permutation of columns.
A destabilization deletes one column label and one row label and changes one value of $O$ or $X$.

\begin{lemma}\label{lem:static}
An exchange of two rows does not change the position of any column, so it does not change the
segments of columns. An exchange of two columns does not change the segments of rows except the
two exchanged columns. Consequently, if a row $r$ is non-interleaving with each of the rows
$s_1,\dots,s_\ell$ that are consecutive to it in this order, then $r$ can be moved past
$s_1,\dots,s_\ell$ by $\ell$ row exchanges, each of which is a grid move. The same holds for columns.
\end{lemma}

\begin{proof}
The first two sentences are immediate from the definitions. The segment of the row $r$ consists of
the positions of the two columns containing its markings, and the segments of $s_1,\dots,s_\ell$
likewise. Row exchanges do not change these positions, so the non-interleaving hypotheses hold at
every step.
\end{proof}

\begin{lemma}[short segments]\label{lem:short}
Let $c$ be a column whose two markings lie in rows with consecutive positions. Then $c$ is
non-interleaving with every column $d$ that has no marking in these two rows. The same holds for a
row whose two markings lie in columns with consecutive positions, and every row with no marking in
these two columns.
\end{lemma}

\begin{proof}
Let the rows of $c$ have positions $x,x+1$. The cyclic interval between them, in either direction,
contains no position or all positions other than $x,x+1$. The markings of $d$ avoid $x,x+1$, so they
lie on one side. In the other direction, $x$ and $x+1$ are both inside or both outside the cyclic
interval between the rows of $d$, because no position lies strictly between them. This is
\lean{columnsNoninterleaving\_of\_adjacent\_rows} and \lean{rowsNoninterleaving\_of\_adjacent\_columns} of
\lean{GridMoveRecognition.lean}, and J10, Lemma~\emph{commutations of a short row or column}.
\end{proof}

\begin{lemma}[separated rows]\label{lem:separated}
Let $W$ be a set of column positions forming an interval $\{p,p+1,\dots,p'\}$ of integers. If the
markings of a row $r$ lie in columns with positions in $W$ and the markings of a row $s$ lie in
columns with positions outside $W$, then $r$ and $s$ are non-interleaving.
\end{lemma}

\begin{proof}
Let $x_0,x_1\in W$ be the positions of the columns of $r$ and $y_0,y_1\notin W$ those of $s$. The
set $(x_0,x_1)^\circ$ is either contained in $W$ or contains every position outside $W$, so $y_0$ and
$y_1$ lie on the same side of it. The set $(y_0,y_1)^\circ$ either contains $W$ or is disjoint from $W$,
because $W$ is an interval of positions avoiding $y_0,y_1$. So $x_0$ and $x_1$ lie on the same side.
\end{proof}

\section{The torus grids}\label{sec:torus}

\begin{definition}\label{def:R}
Let $a\ge2$ and $k\ge1$, and put $N=a+k$. The grid diagram $R_{a,k}$ of size $N$ has
\[
O(c)=-c\bmod N,\qquad X(c)=a-c\bmod N.
\]
\end{definition}

Since $0<a<N$, the residues $-c$ and $a-c$ differ, and both maps are permutations. So
$R_{a,k}$ is a grid diagram.

\begin{lemma}\label{lem:R-reflect}
For $a\ge2$ and $q\ge1$, put $N=a+q$. Then $R_{a,q}$ is obtained from
\[
\lean{reflectedDiagram (torusLink (a-1) (q-1))}
\]
by the cyclic permutation of rows
$\lean{relabelRows (finRotate N)}$. In particular the two diagrams are related by one grid move.
\end{lemma}

\begin{proof}
The diagram $\lean{torusLink (a-1) (q-1)}$ has size $(a-1+1)+(q-1+1)=N$, with $O(c)=c$ and
$X(c)=c+q\bmod N$. Its reflection has $O(c)=N-1-c$ and $X(c)=N-1-(c+q\bmod N)$. Adding one modulo
$N$ to every row gives $O(c)\equiv -c$ and $X(c)\equiv -c-q\equiv a-c$, because $N=a+q$.
\end{proof}

In $R_{a,k}$, the row $t$ has its $O$-marking in column $-t$ and its $X$-marking in column $a-t$,
modulo $N$. We call the rows $1,\dots,a$ the \emph{window rows}. These are exactly the rows whose
$O$-marking lies in one of the last $a$ columns $N-a,\dots,N-1$. The window row $r+1$, for
$0\le r\le a-1$, has its $O$-marking in column $N-1-r$ and its $X$-marking in column $a-1-r$.
We call $r$ its \emph{rank}.

\section{Hybrid grids}\label{sec:hybrid}

Fix $a\ge2$. A positive word $v$ on $a$ strands \emph{covers all ranks} if for every
$r\in\{0,\dots,a-1\}$ some letter of $v$ involves the rank $r$, that is, equals $\sigma_{r-1}$ or
$\sigma_r$. Every word containing all generators covers all ranks, by
\lean{positiveRectilinearWord\_all\_rows\_occur}. We recall the arcs of the admissible grids of a word
$v$ of length $L$ (J10, section \emph{The row-separated word grids}). The corners are the pairs
$(h,c)$ with $0\le h<L$ and $c\in\{0,1,2\}$, in lexicographic order. Every corner is the start of
one horizontal arc, the \emph{arc leaving} the corner. For the letter $\sigma_s$ in position $h$,
the arcs leaving $(h,0),(h,1),(h,2)$ have the heights $s+\frac12$, $s$, $s+1$. The arc leaving
$(h,0)$ ends at $(h,2)$. An arc of integer height $r$ leaving an occurrence of the rank $r$ ends at
the incoming corner of the next occurrence of $r$, where the incoming corner of an occurrence is
$(h,0)$ if $r$ is the lower rank of the letter and $(h,1)$ if it is the upper rank, and the
outgoing corner is $(h,1)$, respectively $(h,2)$.

\begin{definition}\label{def:hybrid}
Let $k\ge1$, $N=a+k$, and let $v$ be a positive word on $a$ strands which is empty or covers all
ranks. The labeled grid diagram $H_{a,k}(v)$ has the following data.
\begin{enumerate}
\item Columns: the labels $\rho_0,\dots,\rho_{N-1}$, followed by the corners of $v$ in
lexicographic order.
\item Rows: the labels $\varrho_t$ for $t\in\{0\}\cup\{a+1,\dots,N-1\}$; for each rank $r$ a label
$\omega_r$ if $v$ is empty and a label $\iota_r$ if $v$ is nonempty; and the arcs of $v$.
\item The $O$-marking of $\rho_c$ is $\varrho_t$ with $t=-c\bmod N$ if this residue is not a window
row, and $\omega_r$, respectively $\iota_r$, if $-c\equiv r+1$ with $0\le r\le a-1$.
\item The $X$-marking of $\rho_c$ is $\varrho_t$ with $t=a-c\bmod N$ if this residue is not a window
row. If $a-c\equiv r+1$ with $0\le r\le a-1$, it is $\omega_r$ when $v$ is empty and the arc of $v$
leaving the outgoing corner of the last occurrence of $r$ in $v$ when $v$ is nonempty.
\item The $O$-marking of a corner of $v$ is the arc leaving it. The $X$-marking of the incoming
corner of the first occurrence of a rank $r$ in $v$ is $\iota_r$. The $X$-marking of every other
corner is the arc of $v$ ending there.
\item Row order: first $\varrho_0$; then the \emph{band}, consisting of the labels
$\omega_r,\iota_r$ of height $r$ and the arcs of $v$ with their heights, ordered by height with an
arbitrary order among equal heights; then $\varrho_{a+1},\dots,\varrho_{N-1}$.
\end{enumerate}
\end{definition}

So $H_{a,k}(v)$ is $R_{a,k}$ with the word grid of $v$ inserted at the right end: the window row of
rank $r$ is cut into the row $\iota_r$, which ends at the first occurrence of $r$ in $v$, and the
arc of $v$ leaving the last occurrence of $r$, which ends in column $\rho_{a-1-r}$.

\begin{lemma}\label{lem:hybrid-grid}
$H_{a,k}(v)$ is a labeled grid diagram. The diagram $H_{a,k}(\varnothing)$ is $R_{a,k}$. Any two
choices of the order among equal heights in the band give diagrams related by row exchanges, each of
which is a grid move.
\end{lemma}

\begin{proof}
Each row label is the $O$-marking of exactly one column: $\varrho_t$ of $\rho_{-t}$, $\omega_r$ and
$\iota_r$ of $\rho_{N-1-r}$, and an arc of $v$ of its starting corner. Each row label is the
$X$-marking of exactly one column: $\varrho_t$ of $\rho_{a-t}$, $\omega_r$ of $\rho_{a-1-r}$,
$\iota_r$ of the first incoming corner of $r$, the arc leaving the last occurrence of $r$ of
$\rho_{a-1-r}$, and every other arc of $v$ of its end corner. For the columns $\rho_c$ the two
markings differ because $-c\not\equiv a-c$, and an $O$-label of the form $\iota_r$ is never an
$X$-label of an $R$-column. For a corner of $v$, the arc leaving it and the label at its $X$-marking
are different, because the matching of $v$ has no fixed point
(\lean{positiveRectilinearWordMatching\_ne\_self}) and because $\iota_r$ is not an arc of $v$.

If $v$ is empty, the rows are $\varrho_0,\omega_0,\dots,\omega_{a-1},\varrho_{a+1},\dots,\varrho_{N-1}$
in this order, and $\omega_r$ occupies the window row $r+1$. So the associated diagram is
$R_{a,k}$.

For the last assertion it suffices, by Lemma~\ref{lem:static} and the argument of J10, Lemma
\emph{same-height rows do not interleave}, to show that two rows of the band of equal height are
non-interleaving. At a half-integer height both rows are middle arcs of different letters of $v$,
whose column intervals are disjoint. At an integer height $r$ the rows are: $\iota_r$, with
columns $\rho_{N-1-r}$ and the first incoming corner $f$ of $r$; the arc $\lambda$ leaving the last
outgoing corner $\ell$ of $r$, with columns $\ell$ and $\rho_{a-1-r}$; and the arcs of $v$ joining
consecutive occurrences of $r$, whose column intervals lie in the closed interval from $f$ to
$\ell$ and are pairwise disjoint (J10, loc.\ cit.). In the linear order of positions,
$\rho_{a-1-r}<\rho_{N-1-r}<f<\ell$. Hence the interval of $\iota_r$ lies inside the interval of
$\lambda$, the arcs between consecutive occurrences lie inside the interval of $\lambda$ and are
disjoint from the interval of $\iota_r$. All endpoints are distinct, because the incoming and
outgoing corners of an occurrence are different. So every pair is nested or disjoint, hence
non-interleaving.
\end{proof}

\section{The descending Coxeter block}\label{sec:block}

Let $\beta=\sigma_{a-2}\sigma_{a-3}\cdots\sigma_0$. Its letter in position $t$ is $\sigma_{a-2-t}$,
for $0\le t\le a-2$. Every rank occurs in $\beta$: the rank $a-1$ only in position $0$, the rank $0$
only in position $a-2$, and a rank $r$ with $1\le r\le a-2$ in positions $a-2-r$ (lower rank) and
$a-1-r$ (upper rank). We use the following names for the corners of $\beta$ and the arcs leaving
them.
\begin{center}
\begin{tabular}{@{}p{0.56\textwidth}p{0.40\textwidth}@{}}
\toprule
$D_t=(t,1)$, $0\le t\le a-2$ & the descending columns\\
$A_r=(a-2-r,0)$, $B_r=(a-2-r,2)$, $0\le r\le a-2$ & the ascending columns of rank $r$\\
$\mu_r$, the arc leaving $A_r$ & height $r+\frac12$, from $A_r$ to $B_r$\\
$e_r$, the arc leaving $D_{a-2-r}$, $1\le r\le a-2$ & height $r$, from $D_{a-2-r}$ to $D_{a-1-r}$\\
$o_0$, the arc leaving $D_{a-2}$; $o_r$, the arc leaving $B_{r-1}$, \mbox{$1\le r\le a-1$} & height $r$,
the last outgoing arc of rank~$r$\\
\bottomrule
\end{tabular}
\end{center}
In a diagram containing $\beta$ as a block of consecutive columns, $\iota_r$ denotes the row of the
$X$-marking of the first incoming corner of $r$, namely $A_r$ for $r\le a-2$ and $D_0$ for $r=a-1$.
Thus the markings of the block columns are
\[
\begin{array}{c|cccc}
 & D_0 & D_t\ (1\le t\le a-3) & D_{a-2}\ (a\ge3) & A_r\\ \hline
 O & e_{a-2} & e_{a-2-t} & o_0 & \mu_r\\
 X & \iota_{a-1} & e_{a-1-t} & e_1 & \iota_r
\end{array}
\qquad
\begin{array}{c|c}
 & B_r\\ \hline
 O & o_{r+1}\\
 X & \mu_r
\end{array}
\]
For $a=2$ the only descending column is $D_0$, with $O=o_0$ and $X=\iota_1$. The column order of the
block is $A_{a-2},D_0,B_{a-2},A_{a-3},D_1,B_{a-3},\dots,A_0,D_{a-2},B_0$.

\begin{definition}\label{def:block-conf}
A labeled grid diagram $G$ is in \emph{block position} if it contains these $3(a-1)$ column labels as
consecutive columns in the order above, with the markings above, where the labels
$\iota_0,\dots,\iota_{a-1},o_0,\dots,o_{a-1},e_1,\dots,e_{a-2},\mu_0,\dots,\mu_{a-2}$ are pairwise
distinct, except that $\iota_r=o_r$ is allowed, and if the row order satisfies:
\begin{enumerate}
\item[(B1)] for every $r$, all rows among $\iota_r,o_r,e_r$ lie below $\mu_r$, and $\mu_r$ lies below
all rows among $\iota_{r+1},o_{r+1},e_{r+1}$;
\item[(B2)] $e_r$ lies below $\iota_r$ for $1\le r\le a-2$, and $o_0$ lies below $\iota_0$ unless
$o_0=\iota_0$.
\end{enumerate}
\end{definition}

The next lemma performs the two first stages of the reduction. Its moves only use the rows and
columns of the block, together with Lemma~\ref{lem:short}, so the rest of the diagram is arbitrary.

\begin{lemma}[block reduction]\label{lem:block}
Let $G$ be in block position, and assume $\iota_r\ne o_r$ for every $r$. Then $G\Mov G'$, where
$G'$ is obtained from $G$ as follows: the block columns are replaced by the $a$ consecutive columns
\[
V,\ B_{a-2},\ B_{a-3},\ \dots,\ B_0,
\]
with $O(V)=o_0$, $X(V)=\iota_{a-1}$, $O(B_r)=o_{r+1}$ and $X(B_r)=\iota_r$; the rows
$e_1,\dots,e_{a-2},\mu_0,\dots,\mu_{a-2}$ are deleted; and the relative order of all other rows,
and of all columns outside the block, is unchanged.
\end{lemma}

\begin{proof}
\emph{Stage 1: merging the descending columns.} We show by induction on $t\in\{0,\dots,a-2\}$ that
$G$ is related to the diagram $G_t$ in which the block columns are
\[
A_{a-2},\ M_t,\ B_{a-2},A_{a-3},B_{a-3},\dots,A_{a-2-t},B_{a-2-t},\ A_{a-3-t},D_{t+1},B_{a-3-t},\dots,
A_0,D_{a-2},B_0,
\]
where $M_t$ is the label $D_t$ with $X(M_t)=\iota_{a-1}$ and the other markings unchanged, and the
rows $e_{a-1-t},\dots,e_{a-2}$ (an empty list for $t=0$) are deleted. For $t=0$ this is $G$. The
case $t=a-2$ gives the order $A_{a-2},V,B_{a-2},A_{a-3},B_{a-3},\dots,A_0,B_0$ with $V=M_{a-2}$,
$O(V)=o_0$, $X(V)=\iota_{a-1}$.

Let $t\le a-3$. The column $D_{t+1}$ has $X$-marking $e_{a-2-t}$ and $O$-marking $e_{a-3-t}$ if
$t\le a-4$, and $o_0$ if $t=a-3$. Its rows have heights $a-2-t$ and $a-3-t$. We move $D_{t+1}$ to the
left until it is immediately right of $M_t$. It crosses the columns $A_{a-3-t}$ and $B_s,A_s$ for
$a-2-t\le s\le a-2$, $s\ne a-2$ for $A_s$. By Lemma~\ref{lem:static} it suffices to check that each of them is
non-interleaving with $D_{t+1}$, and this depends only on the order of their rows.
\begin{itemize}
\item $A_{a-3-t}$ has rows $\iota_{a-3-t}$ and $\mu_{a-3-t}$. By (B2) the $O$-row of $D_{t+1}$ lies
below $\iota_{a-3-t}$. By (B1) $\mu_{a-3-t}$ lies below $e_{a-2-t}$. So the segment of $A_{a-3-t}$
lies strictly inside the segment of $D_{t+1}$ in the linear order of rows. Nested segments with
distinct endpoints are non-interleaving.
\item $A_s$ with $a-2-t\le s\le a-3$ has rows $\iota_s,\mu_s$ of heights $s,s+\frac12$. For
$s=a-2-t$, (B2) puts $e_{a-2-t}$ below $\iota_s$, and for larger $s$ the heights separate. So the
segments are disjoint.
\item $B_s$ with $s\ge a-2-t$ has rows $\mu_s,o_{s+1}$ of heights $s+\frac12,s+1$, all above the rows
of $D_{t+1}$ by (B1). The segments are disjoint.
\end{itemize}
Now $M_t$ and $D_{t+1}$ are consecutive columns and the row $e_{a-2-t}$ has its $O$-marking in $M_t$
and its $X$-marking in $D_{t+1}$. Move this row up until it is immediately below $\iota_{a-1}$. The
only rows with markings in $M_t$ or $D_{t+1}$ are $\iota_{a-1}$, $e_{a-2-t}$ and the $O$-row of
$D_{t+1}$, which lies below $e_{a-2-t}$. So every row crossed has no marking in these two columns,
and Lemma~\ref{lem:short} applies at every step. Now the corners $(M_t,e_{a-2-t})$,
$(D_{t+1},e_{a-2-t})$, $(M_t,\iota_{a-1})$ carry $O,X,X$, and $(D_{t+1},\iota_{a-1})$ carries no
$O$, because the $O$-row of $D_{t+1}$ has height at most $a-3-t<a-1$. The $X$-destabilization
deletes $M_t$ and $e_{a-2-t}$ and gives $X(D_{t+1})=\iota_{a-1}$. Renaming $D_{t+1}$ as $M_{t+1}$
gives $G_{t+1}$. The rows other than $e_{a-2-t}$ keep their relative order.

\emph{Stage 2: removing the middle arcs.} Treat $r=a-2,a-3,\dots,0$ in turn. If $r=a-2$, first move
$A_{a-2}$ to the right past $V$. Its segment, with rows $\iota_{a-2}$ and $\mu_{a-2}$, lies strictly
inside the segment of $V$, with rows $o_0$ and $\iota_{a-1}$: by (B1) and (B2), $o_0$ lies below
$\iota_{a-2}$ (for $a=2$ this is $o_0<\iota_0$) and $\mu_{a-2}$ lies below $\iota_{a-1}$. For $r\le
a-3$ the columns $A_r$ and $B_r$ are already consecutive. Now the row $\mu_r$ has its markings in
the consecutive columns $A_r,B_r$. Move it down until it is immediately above $\iota_r$. The rows
with markings in $A_r$ or $B_r$ are $\iota_r$, $\mu_r$ and $o_{r+1}$, and $o_{r+1}$ lies above
$\mu_r$, so Lemma~\ref{lem:short} applies at every step. The corners $(A_r,\mu_r)$, $(B_r,\mu_r)$,
$(A_r,\iota_r)$ carry $O,X,X$ and $(B_r,\iota_r)$ carries no $O$, because $O(B_r)=o_{r+1}$. The
$X$-destabilization deletes $A_r$ and $\mu_r$ and gives $X(B_r)=\iota_r$.

After both stages the block columns are $V,B_{a-2},\dots,B_0$ with the stated markings. Each step
deleted one of the listed rows or exchanged rows only to bring a row that is later deleted next to
another row, so the relative order of the remaining rows is unchanged.
\end{proof}

\section{Absorption of one block}\label{sec:absorb}

\begin{lemma}[absorption]\label{lem:absorb}
Let $a\ge2$, $k\ge1$, and let $v$ be a positive word on $a$ strands which is empty or covers all
ranks. Then $H_{a,k}(\beta v)\Mov H_{a,k+1}(v)$.
\end{lemma}

\begin{proof}
Put $N=a+k$. The word $\beta v$ covers all ranks, so $H=H_{a,k}(\beta v)$ has the rows $\iota_r$. In
$\beta v$ the first occurrence of every rank lies in $\beta$, so $\iota_r$ is the $X$-marking of the
first incoming corner of $r$ in the block $\beta$, and the notation of Section~\ref{sec:block}
applies. The arc $o_r$ leaving the last occurrence of $r$ in $\beta$ ends at the first occurrence of
$r$ in $v$ if $v$ is nonempty, and in the column $\rho_{a-1-r}$ if $v$ is empty. By
Lemma~\ref{lem:hybrid-grid} we may choose the order among rows of equal height in the band. We
choose it so that, for every $r$, the row $o_r$ is the lowest and $\iota_r$ the highest row of height
$r$. Then (B1) holds because the band is ordered by height, and (B2) holds by this choice. Also
$\iota_r\ne o_r$, since $\iota_r$ is not an arc. So $H$ is in block position and
Lemma~\ref{lem:block} gives $H\Mov H'$, where the block is replaced by $V,B_{a-2},\dots,B_0$.

\emph{Stage 3.} For $r=0,1,\dots,a-2$ in turn we do the following. The column $\rho_{N-1-r}$ has
$O$-marking $\iota_r$, and $X(B_r)=\iota_r$, $O(B_r)=o_{r+1}$.
\begin{enumerate}
\item Move the row $\iota_r$ up until it is immediately below $o_{r+1}$. In the current diagram,
the markings of $\iota_r$ lie in $\rho_{N-1-r}$ and $B_r$. Since $\iota_r$ is the highest row of height
$r$ and $o_{r+1}$ is the lowest row of height $r+1$, and since the middle arcs of the block were
deleted in Lemma~\ref{lem:block}, every row strictly between them is a middle arc of $v$ of height
$r+\frac12$. Such an arc has both markings in columns of $v$. The columns $\rho_0,\dots,\rho_{N-1}$
followed by the block columns form an interval of positions containing both markings of
$\iota_r$, and the columns of $v$ lie outside it. By Lemma~\ref{lem:separated} and
Lemma~\ref{lem:static}, all these exchanges are grid moves.
\item Move the column $B_r$ to the left until it is immediately right of $\rho_{N-1-r}$. The rows of
$B_r$ now have consecutive positions. The columns crossed are $V$, $B_{a-2},\dots,B_{r+1}$ and
$\rho_{N-r},\dots,\rho_{N-1}$. Their markings are: $o_0,\iota_{a-1}$ for $V$; $o_{s+1},\iota_s$ for
$B_s$; and for $\rho_{N-1-s}$ with $s<r$, the $O$-marking $o_{s+1}$ (it was changed in the step for
$s$) and an $X$-marking. This $X$-marking is the label of the residue $a-(N-1-s)\equiv s+1-k$
modulo $N$: a row $\varrho_t$ if $s<k$, and otherwise the window row of rank $u=s-k$, namely an arc
of $v$, or the row $o_u$ if $v$ is empty. Since $u\le s<r$, none of these markings lies in the rows
$\iota_r$, $o_{r+1}$. So Lemma~\ref{lem:short} applies at every step.
\item The corners $(\rho_{N-1-r},\iota_r)$, $(B_r,\iota_r)$, $(B_r,o_{r+1})$ carry $O,X,O$, and the
corner $(\rho_{N-1-r},o_{r+1})$ carries no $X$-marking, by the description of $X(\rho_{N-1-r})$ just
given with $s=r$ (for $v$ empty and $r\ge k$, $X(\rho_{N-1-r})=o_{r-k}\ne o_{r+1}$). The
$O$-destabilization deletes $B_r$ and $\iota_r$ and gives $O(\rho_{N-1-r})=o_{r+1}$.
\end{enumerate}

\emph{Identification.} After stage 3 the columns are $\rho_0,\dots,\rho_{N-1},V$, followed by the
corners of $v$. The markings are those of $H_{a,k}(\beta v)$ for all columns $\rho_c$ except
$O(\rho_{N-1-r})=o_{r+1}$ for $0\le r\le a-2$; moreover $O(V)=o_0$, $X(V)=\iota_{a-1}$, and the corners
of $v$ keep their markings, with $X$-marking $o_r$ at the first incoming corner of $r$ in $v$.
Rename the labels by
\[
\rho_c\mapsto\rho'_c\ (c<N),\quad V\mapsto\rho'_N,\quad \varrho_0\mapsto\varrho'_0,\quad
\varrho_t\mapsto\varrho'_{t+1}\ (t\ge a+1),\quad \iota_{a-1}\mapsto\varrho'_{a+1},
\]
\[
o_r\mapsto\iota'_r\ \text{if $v$ is nonempty},\qquad o_r\mapsto\omega'_r\ \text{if $v$ is empty},
\]
and keep the arcs and corners of $v$. We compare with $H_{a,k+1}(v)$, of size $N+1$, column by
column, using residues modulo $N+1$ for $H_{a,k+1}(v)$ and modulo $N$ for $H_{a,k}(\beta v)$.
\begin{itemize}
\item $O$-markings. In $H_{a,k}(\beta v)$ after stage 3: $O(\rho_0)=\varrho_0$;
$O(\rho_c)=\varrho_{N-c}$ for $1\le c\le k-1$; $O(\rho_k)=\iota_{a-1}$; $O(\rho_c)=o_{N-c}$ for
$k+1\le c\le N-1$; and $O(V)=o_0$. In $H_{a,k+1}(v)$: $O(\rho'_0)=\varrho'_0$;
$O(\rho'_c)=\varrho'_{N+1-c}$ for $1\le c\le k$, since then $N+1-c\ge a+1$; and $O(\rho'_c)$ is the
row $\iota'_{N-c}$ (or $\omega'_{N-c}$) for $k+1\le c\le N$. Under the renaming these agree: for
$c\le k-1$ because $\varrho_{N-c}\mapsto\varrho'_{N-c+1}$, for $c=k$ because
$\iota_{a-1}\mapsto\varrho'_{a+1}=\varrho'_{N+1-k}$, and for $k+1\le c\le N$ because
$o_{N-c}\mapsto\iota'_{N-c}$ (with $V=\rho'_N$ and $o_0\mapsto\iota'_0$).
\item $X$-markings of the columns $\rho_c$, $c\le N-1$. For $0\le c\le a-1$ both diagrams have the
window row of rank $a-1-c$, which is the arc of $v$ leaving its last occurrence of $a-1-c$ if $v$ is
nonempty, and $o_{a-1-c}\mapsto\omega'_{a-1-c}$ if $v$ is empty. For $c=a$ both have the row of
index $0$, that is $\varrho_0\mapsto\varrho'_0$. For $a+1\le c\le N-1$, the diagram
$H_{a,k}(\beta v)$ has $\varrho_{N+a-c}$ and $H_{a,k+1}(v)$ has $\varrho'_{N+1+a-c}$, which agree.
\item $X(V)=\iota_{a-1}\mapsto\varrho'_{a+1}$, and in $H_{a,k+1}(v)$ the column $\rho'_N$ has
$X$-marking $a-N\equiv a+1\pmod{N+1}$; since $k\ge1$, $a+1\le N$ is not a window row.
\item Corners of $v$: the incoming corner of the first occurrence of $r$ has $X$-marking $o_r$,
which is $\iota'_r$; the other markings are arcs of $v$ in both diagrams. When $v$ is empty there
are no such columns.
\end{itemize}
The rows of the renamed diagram are $\varrho'_0$, then the band rows, then $\varrho'_{a+1}=\iota_{a-1}$
and $\varrho'_{a+2},\dots,\varrho'_{N}$. Indeed $\iota_{a-1}$ was the highest row of height $a-1$ and
nothing above it in the band was moved, and the row exchanges of stage 3 only moved rows $\iota_r$
inside the band before deleting them. The band rows are ordered by height, because each $o_r$ kept
its position, the deleted rows were the only ones moved, and the arcs of $v$ never moved. So the
renamed diagram is $H_{a,k+1}(v)$ with some order among equal heights, and
Lemma~\ref{lem:hybrid-grid} completes the proof.
\end{proof}

\section{The first block}\label{sec:base}

\begin{lemma}[base case]\label{lem:base}
Let $a\ge2$ and let $v$ be a nonempty positive word on $a$ strands that covers all ranks. Then every
admissible grid of $\beta v$ is related to $H_{a,1}(v)$.
\end{lemma}

\begin{proof}
By \lean{IsAdmissibleWordGrid.movesTo} we may choose the admissible grid. View it as the labeled
diagram whose columns are the corners of $\beta v$ and whose rows are its arcs, ordered by height with
the following choice among equal heights: $o_r$ lowest and the arc $\iota_r$ highest among the arcs
of height $r$. Here $\iota_r$ is the arc of $v$ leaving the last occurrence of $r$ in $v$; it ends at
the first occurrence of $r$, which lies in $\beta$. Since $v$ covers all ranks, $\iota_r\ne o_r$, and
the diagram is in block position. Lemma~\ref{lem:block} replaces the block by $V,B_{a-2},\dots,B_0$.

Perform the $O$-stabilization of the $O$-marking of $V$ that inserts a new column $\kappa$
immediately right of $V$ and a new row $\nu$ immediately below $o_0$: afterwards $O(V)=\nu$,
$O(\kappa)=o_0$ and $X(\kappa)=\nu$. This is the TauCeti construction \lean{stabilizeO} with the new
column and the new row cyclically adjacent to the split column and to its $O$-row, so it is an
elementary stabilization. Next move $\nu$ down to the bottom. Its markings lie in the consecutive
columns $V,\kappa$, and the rows with markings in these columns are $\nu$, $o_0$ and $\iota_{a-1}$,
both above $\nu$. So Lemma~\ref{lem:short} applies. Next move $\kappa$ to the right past
$B_{a-2},\dots,B_0$. The segment of $\kappa$ has rows $\nu$ and $o_0$, and the segment of $B_r$ has
rows $\iota_r$ and $o_{r+1}$, which lie above $o_0$ (for $r=0$ because $o_0$ is the lowest row of
height $0$). The segments are disjoint, so these are grid moves.

The columns are now $V,B_{a-2},\dots,B_0,\kappa$ followed by the corners of $v$, and the rows are
$\nu$ followed by the arcs of $\beta v$ that were not deleted, in an order by height. Rename
$V\mapsto\rho_0$, $B_r\mapsto\rho_{a-1-r}$, $\kappa\mapsto\rho_a$, $\nu\mapsto\varrho_0$ and
$o_r\mapsto\iota'_r$, keeping the arcs and corners of $v$. In $H_{a,1}(v)$, of size $a+1$, the column
$\rho_c$ has $O$-marking $-c\bmod(a+1)$ and $X$-marking $a-c$. Thus $\rho_0$ has $O=\varrho_0$ and $X$
the window row of rank $a-1$, that is the arc of $v$ leaving its last occurrence of $a-1$; this arc is
$\iota_{a-1}=X(V)$. For $1\le c\le a-1$, $\rho_c=B_{a-1-c}$ has $O$ the window row $a+1-c$ of rank
$a-c$, which is $\iota'_{a-c}=o_{a-c}=O(B_{a-1-c})$, and $X$ the window row of rank $a-1-c$, which is the
arc $\iota_{a-1-c}=X(B_{a-1-c})$. Finally $\rho_a=\kappa$ has $O=\iota'_0=o_0$ and $X=\varrho_0=\nu$.
The corners of $v$ have the markings of Definition~\ref{def:hybrid}, with $X$-marking $o_r=\iota'_r$ at
the first incoming corner of $r$ in $v$. The rows are $\varrho_0$ followed by the band, ordered by height.
\end{proof}

\section{Proof of Theorem~\ref{thm:main}}\label{sec:proof}

\subsection{Conjugacy of the words}
\begin{lemma}\label{lem:conj-power}
Let $\gamma,\gamma'$ be positive words related by far commutations inside the word and by rotations
$u\,i\mapsto i\,u$ (the relation \lean{FarRotationEquiv}), and let $q\ge1$. If $\gamma$ is nonempty, then
$\gamma^q$ and $\gamma'^q$ are related by \lean{PositiveWordConj}.
\end{lemma}

\begin{proof}
It suffices to treat one step. If $\gamma=u\,i\,j\,v\mapsto u\,j\,i\,v$ is a far commutation, then
$\gamma^q\mapsto\gamma'^q$ is the composition of $q$ far commutations, one in each factor. If
$\gamma=u\,i$ and $\gamma'=i\,u$, then $\gamma^q=(u\,i)^q$, and the rotation moving the last letter
$i$ to the front gives $i\,(u\,i)^{q-1}u=(i\,u)^q$. Both are generators of \lean{PositiveWordConj}. The
words in each step are nonempty, so reflexivity is never needed.
\end{proof}

The word $\beta$ is a permutation of $\delta=\lean{List.finRange (a-1)}$, so
$\lean{farRotationEquiv\_coxeterWord}$ gives $\beta\approx\delta$ in \lean{FarRotationEquiv}.
For $a\ge2$ both words are nonempty, and Lemma~\ref{lem:conj-power} gives
$\beta^q\sim\delta^q$ in \lean{PositiveWordConj}. Every generator occurs in $\delta^q$.

\begin{lemma}[cyclic and far moves, J09]\label{lem:J09}
Let $u,v$ be positive words such that $v$ is obtained from $u$ by one rotation or one far commutation.
Then $\lean{wordGrid u}\Mov\lean{wordGrid v}$.
\end{lemma}

\begin{proof}
This is proved in \lean{docs/route-n/J09-word-grid-moves.tex}, Propositions on rotations and far
commutations. For a rotation $u=u'\,i\mapsto i\,u'$, three cyclic permutations of columns move the
three columns of the last letter to the front, and the result is an admissible grid of $i\,u'$. For a
far commutation $u'\,i\,j\,v'\mapsto u'\,j\,i\,v'$ with $|i-j|\ge2$, the rows of the columns of the
letter $i$ have heights in $[i,i+1]$ and those of the letter $j$ in $[j,j+1]$; these intervals are
disjoint, so the nine column exchanges that move the three columns of $j$ before those of $i$ are
grid moves, and the result is an admissible grid of $u'\,j\,i\,v'$. In both cases
\lean{IsAdmissibleWordGrid.movesTo} relates the result to the selected grid.
\end{proof}

Applying Lemma~\ref{lem:J09} along a sequence of rotations and far commutations gives
\begin{equation}\label{eq:conj}
\lean{wordGrid}(\delta^q)\Mov\lean{wordGrid}(\beta^q).
\end{equation}

\subsection{The case \texorpdfstring{$q\ge2$}{q at least 2}}
Let $q\ge2$. The word $\beta^{q-1}$ is nonempty and covers all ranks. Lemma~\ref{lem:base} gives
$\lean{wordGrid}(\beta^q)\Mov H_{a,1}(\beta^{q-1})$. For $1\le k\le q-1$, Lemma~\ref{lem:absorb} with
$v=\beta^{q-1-k}$ gives $H_{a,k}(\beta^{q-k})\Mov H_{a,k+1}(\beta^{q-1-k})$. The word
$\beta^{q-1-k}$ is empty or covers all ranks. Composing,
$\lean{wordGrid}(\beta^q)\Mov H_{a,q}(\varnothing)=R_{a,q}$ by Lemma~\ref{lem:hybrid-grid}. With
\eqref{eq:conj} and Lemma~\ref{lem:R-reflect} this proves the theorem for $q\ge2$.

\subsection{The case \texorpdfstring{$q=1$}{q equal to 1}}
For $q=1$ the right side is the reflection of $\lean{torusLink (a-1) 0}$, which is related by
Lemma~\ref{lem:R-reflect} to $R_{a,1}$. In $R_{a,1}$, of size $a+1$, every column has
$X(c)\equiv O(c)-1$. We show that $R_{a,1}$ and the word grid of $\delta$ are both related to the
diagram $U=(O,X)=((0,1),(1,0))$ of size $2$.

\emph{The diagram $R_{a,1}$.} Let $n\ge3$ and let $G$ be a grid diagram of size $n$ with $O(c)=-c$
and $X(c)=-c-1$ modulo $n$. The columns $c=n-2$ and $c+1=n-1$ are consecutive; the rows $2$ and $1$
are consecutive; the corners $(n-2,2),(n-1,1),(n-2,1)$ carry $O,O,X$, since $O(n-2)=2$, $O(n-1)=1$,
$X(n-2)=1$, and $(n-1,2)$ carries no $X$ because $X(n-1)=0$. So this is an $O$-destabilization deleting
column $n-2$ and row $1$ and setting $O(n-1)=2$. After renumbering, the remaining diagram has
$O(c)=-c$, $X(c)=-c-1$ modulo $n-1$: for $c\le n-3$ the rows $-c,-c-1$ are at least $2$ or equal to
$0$, and they are renumbered by subtracting one when they exceed $1$, which turns residues modulo $n$
into residues modulo $n-1$; the last column $n-2$ has $O=2-1=1\equiv-(n-2)$ and $X=0\equiv -(n-2)-1$.
By induction, $R_{a,1}\Mov G_2$, where $G_2$ has $O=(0,1)$ and $X=(1,0)$, which is $U$.

\emph{The word grid of $\delta$.} Let $w=\delta$ on $a$ strands, of length $a-1$. In an admissible grid,
the rank $0$ occurs only in letter $0$, so the arc leaving $(0,1)$ is the only arc of height $0$, it
ends at $(0,0)$, and it is the lowest row. The arc leaving $(0,0)$ has height $\frac12$ and is the
second row. So column $0$ has $X$ in row $0$ and $O$ in row $1$, column $1$ has $O$ in row $0$, and
$X(1)$ is the arc ending at the incoming upper corner of $\sigma_0$, of height $1$, so $X(1)\ne 1$. The
corners $(0,0),(0,1),(1,0)$ carry $X,O,O$, and $(1,1)$ carries no $X$: this is an $O$-destabilization
deleting column $0$ and row $0$, with $O(1)=1$. We show by induction on $a$ that the word grid of
$\delta$ on $a$ strands is related to $U$; for $a=2$ the word is $\sigma_0$, of size $3$, and the
destabilization just described gives a diagram of size $2$, which is $U$ or its cyclic row shift.
For $a\ge3$ we perform, after the destabilization above, two further destabilizations which remove
the remaining columns $(0,1)$ and $(0,2)$ of the first letter, and we verify that the result is an
admissible grid of $\sigma_1\cdots\sigma_{a-2}$ on the strands $1,\dots,a-1$, which is the word grid
of $\delta$ on $a-1$ strands after relabeling ranks.

Write the arcs of $\delta$ (for $q=1$) as follows: $\mu_h$ leaving $(h,0)$; for $1\le r\le a-2$, the arc
$c_r$ leaving $(r-1,2)$ and ending at $(r,0)$, and the arc $d_r$ leaving $(r,1)$ and ending at
$(r-1,1)$; the arc $d_0$ leaving $(0,1)$ and ending at $(0,0)$; and the arc $c_{a-1}$ leaving
$(a-2,2)$ and ending at $(a-2,1)$. Choose the order among equal heights with $c_r$ below $d_r$. The rows
are then
\[
d_0<\mu_0<c_1<d_1<\mu_1<c_2<d_2<\cdots<\mu_{a-2}<c_{a-1}.
\]
After the first destabilization (deleting column $(0,0)$ and row $d_0$, with $O((0,1))=\mu_0$), the
columns $(0,1)$ and $(0,2)$ are the first two columns, with markings $O=\mu_0,X=d_1$ and
$O=c_1,X=\mu_0$. The rows $\mu_0$ and $c_1$ are the two lowest rows. The corners
$((0,2),c_1),((0,2),\mu_0),((0,1),\mu_0)$ carry $O,X,O$ and $((0,1),c_1)$ carries no $X$, since
$X((0,1))=d_1\neq c_1$. The $O$-destabilization deletes column $(0,2)$ and row $\mu_0$ and gives
$O((0,1))=c_1$. Now column $(0,1)$ is first, with $O=c_1$ and $X=d_1$, the rows $c_1,d_1$ are the two
lowest rows, and column $(1,0)$ is second, with $X=c_1$ and $O=\mu_1$. The corners
$((0,1),c_1),((0,1),d_1),((1,0),c_1)$ carry $O,X,X$ and $((1,0),d_1)$ carries no $O$. The
$X$-destabilization deletes column $(0,1)$ and row $c_1$ and gives $X((1,0))=d_1$. In the remaining
diagram the arc $d_1$ leaves $(1,1)$ and ends at $(1,0)$, and it is the lowest row; all other markings
are those of the letters $1,\dots,a-2$. This is the admissible grid of $\sigma_1\cdots\sigma_{a-2}$ in
which $d_1$ plays the role of $d_0$, with the same order convention. By induction the word grid of
$\delta$ is related to $U$ or to its cyclic row shift, and $U$ is related to its cyclic row shift by
one cyclic permutation. Hence $\lean{wordGrid}(\delta)\Mov R_{a,1}$, and the theorem holds for $q=1$.

\begin{remark}
The three destabilizations per letter in the case $q=1$ were found by machine and agree with the
sequence printed by the search: an $O$-destabilization, then an $X$-destabilization, at the first
columns. The verification above is by hand.
\end{remark}

\section{Formalization}\label{sec:formalization}
The case $q\ge2$ of Theorem~\ref{thm:main} is formalized without additional hypotheses in
\lean{TorusWordGridMoves.lean} (\lean{wordGrid\_torusWord\_movesTo} and, for $a=S$ and $q=S+1$,
\lean{reflectedTorusGrid\_movesTo\_coxeterPower}). The case $q=1$ is not formalized. The formal
proof follows Sections~\ref{sec:hybrid}--\ref{sec:proof} with the following changes of
bookkeeping.
\begin{enumerate}
\item Lemma~\ref{lem:hybrid-grid} allows any order among rows of equal height. The formalization
fixes one order (\lean{hybLtS}): at each integer height the arcs of later blocks lie lower, the
arc of the last block leaving its last occurrence of the rank lies above all other arcs, and
$\iota_r$ lies on top. With this order the hypotheses of Lemma~\ref{lem:block} hold.
\item After stage~3 of Lemma~\ref{lem:absorb}, and after the moves of Lemma~\ref{lem:base}, the
arc $o_r$ of the first block lies directly below the arc of the last block leaving its last
occurrence of $r$. The column interval of the first arc lies strictly inside the column interval
of the second, so one commutation exchanges the two rows (\lean{LabeledGrid.exists\_swapOuts}).
These exchanges restore the fixed order and replace the appeal to
Lemma~\ref{lem:hybrid-grid}.
\item The formal block reduction (\lean{exists\_blockReduction}) leaves the columns in the order
$B_{a-2},\dots,B_1,V,B_0$. In Lemma~\ref{lem:absorb} the column $V$ is first moved to the right
past $B_0$. In Lemma~\ref{lem:base} the $O$-stabilization is performed at $V$ in this position,
$\nu$ is moved to the bottom, $\kappa$ is moved to the right past $B_0$, and then $V$ is moved to
the left past $B_{a-2},\dots,B_1$. Each of these exchanges is a commutation of columns whose
segments are nested or disjoint.
\item The admissible grid of $\beta^q$ used in Lemma~\ref{lem:base} is constructed with the cyclic
matching of $\beta^q$ computed explicitly (\lean{betaPow\_matching\_one},
\lean{betaPow\_matching\_two}), and it is related to $\lean{wordGrid}(\beta^q)$ by
\lean{IsAdmissibleWordGrid.movesTo}.
\end{enumerate}

\section{Summary of dependencies}
\begin{longtable}{p{0.48\textwidth}p{0.44\textwidth}}
\toprule
Ingredient & Source\\
\midrule
Grid moves, \lean{MovesTo} & \lean{External/TauCeti/KnotTheory/Grid/Move/Basic.lean}\\
Destabilization from a $2\times2$ pattern & \lean{isXStabilization\_eraseOCorner},
\lean{isDestabilization\_swapMarkings}\\
Short segments & \lean{columnsNoninterleaving\_of\_adjacent\_rows},
\lean{rowsNoninterleaving\_of\_adjacent\_columns}\\
Admissible grids, same-height rows & \lean{PositiveWordGridAdmissible.lean}, J10 chapter\\
$\beta\approx\delta$ & \lean{farRotationEquiv\_coxeterWord} (J08)\\
Rotations and far commutations of word grids & Lemma~\ref{lem:J09}, J09\\
Lemmas~\ref{lem:static}--\ref{lem:separated}, \ref{lem:hybrid-grid}--\ref{lem:base} & new, proved here\\
\bottomrule
\end{longtable}

No step uses a geometric isotopy, an identification of grid diagrams with links, or a numerical
computation. The chirality computation of Section~\ref{sec:chirality} is evidence only.

\end{document}
